% !TeX root = ../thuthesis-example.tex

\begin{resume}

  \section*{个人简历}

  1999 年 12 月出生于重庆市石柱土家族自治县。

  2018 年 9 月考入北京航空航天大学生物与医学工程学院，2022 年 6 月本科毕业并获得工学学士学位。

  2023 年 9 月考入清华大学深圳国际研究生院攻读电子信息硕士至今。


  \section*{在学期间完成的相关学术成果}

  \subsection*{学术论文}

  \begin{achievements}
    \item Jiahui Liu$^{\#}$, Yun Chen$^{\#}$, \textbf{Jian Tang}, Xupu Xing, Jinshun Lin, Juping Sun, Xin-Hui Xing, Juan Li, Can Yang Zhang$^{*}$. Deep Learning-Driven Integrated Pipeline for De Novo Design and Synthesis of Antimicrobial Peptides. In press[J]. (已被npj Drug Discovery接收)
  \end{achievements}

\end{resume}



% 本科生格式：

% \begin{resume}
%   \section*{学术论文}
%
%   \begin{achievements}
%     \item ZHOU R, HU C, OU T, et al. Intelligent GRU-RIC Position-Loop
%       Feedforward Compensation Control Method with Application to an
%       Ultraprecision Motion Stage[J], IEEE Transactions on Industrial
%       Informatics, 2024, 20(4): 5609-5621.
%
%     \item 杨轶, 张宁欣, 任天令, 等. 硅基铁电微声学器件中薄膜残余应力的研究[J].
%       中国机械工程, 2005, 16(14):1289-1291.
%
%     \item YANG Y, REN T L, ZHU Y P, et al. PMUTs for handwriting recognition.
%       In press[J]. (已被Integrated Ferroelectrics录用)
%
%   \end{achievements}
%
%
%   \section*{专利}
%
%   \begin{achievements}
%     \item 胡楚雄, 付宏, 朱煜, 等. 一种磁悬浮平面电机: ZL202011322520.6[P]. 2022-04-01.
%
%     \item REN T L, YANG Y, ZHU Y P, et al. Piezoelectric micro acoustic sensor
%       based on ferroelectric materials: No.11/215, 102[P]. (美国发明专利申请号.)
%
%   \end{achievements}
% \end{resume}
